\section{Transformer 为什么有效：表达、优化、硬件与上下文学习}

前面已经说明 Transformer “能做什么”。本章回答更难的问题：为什么这套 package 在规模化后不断获得新能力。Karpathy 的答案不是 attention 单因素，而是三种性质同时成立：forward pass 足够 expressive，residual/normalization 让梯度优化可行，shallow-wide graph 又能高效使用 GPU；大规模 next-token training 最终产生 in-context adaptation。

\lecturefigure{slide-52-what-makes-transformer-effective.jpg}{核心追问：到底什么使 Transformer 如此有效？}{Stanford Online 官方视频 00:55:39--00:55:48。}

这张过渡页应带着四个尺度阅读：单层尺度上，attention 路由信息；深度尺度上，residual blocks 累积计算；训练尺度上，数据、参数与 compute 共同扩展；推理尺度上，prompt 改变 activations，使同一组 weights 表现出不同任务行为。任何只在一个尺度解释成功的理论都不完整。

\subsection{GPT-3：few-shot prompt 暴露 in-context learning}

GPT-3 论文将任务说明和若干 input--output examples 放进 prompt，不更新参数也能提高新样例准确率。\term{few-shot learning} 在这里指 context 内给少量示例；\term{zero-shot} 指只给任务描述；\term{fine-tuning} 则真正用 gradient descent 更新 weights。三者不能混用。

\lecturefigure{slide-53-gpt3.jpg}{GPT-3：规模化自回归语言模型与 few-shot evaluation。}{Stanford Online 官方视频 00:55:43--00:56:31。}

\begin{knowledgebox}{读图：论文标题低估了机制问题}
“Language Models are Few-Shot Learners”报告的是行为现象：更多 context examples 往往提高任务表现。Karpathy 更愿称之为 in-context learning 或 meta-learning，因为模型似乎在一次 forward pass 的 activations 中适配任务。现象不等于机制；具体任务还会受示例顺序、格式、label words 与数据污染影响。
\end{knowledgebox}

\lecturefigure{slide-54-chatgpt.jpg}{ChatGPT 作为 2023 年初的公共交互案例。}{Stanford Online 官方视频 00:55:43--00:56:31。}

课程把 ChatGPT 放在 GPT-3 之后，是为了说明 prompt-conditioned behavior 已进入大众产品，而不是展开其训练细节。Karpathy 在 Q\&A 明确说自己不知道 ChatGPT internals。讲义因此只保留可观察层：用户通过自然语言 context 指定任务，模型在固定权重下产生条件化输出；不臆测未公开参数或系统实现。

\subsection{Scaling laws：性能改善是经验规律，不是无限承诺}

Scaling-law 研究用经验函数描述 loss 随参数 $N$、数据量 $D$ 或 compute $C$ 的变化，例如示意性写成

\[
L(N)\approx L_\infty+aN^{-\alpha},\qquad
L(D)\approx L_\infty+bD^{-\beta}.
\]

其中，$L_\infty$ 是不可约损失近似，$a,b,\alpha,\beta$ 由特定模型族与数据拟合。幂律支持“可预测地扩展”这一工程策略，却不保证所有能力单调、所有数据同质或成本永远合理。

\lecturefigure{slide-55-scaling-laws.jpg}{Scaling laws：模型、数据与计算扩展带来可测的 loss 趋势。}{Stanford Online 官方视频 00:56:31--00:57:14。}

\begin{warningbox}{读图：曲线外推有适用域}
先看横轴是 parameters、tokens 还是 compute，再看纵轴是训练/验证 loss 还是具体 downstream metric；不同曲线不能直接互换。幂律来自观察区间内的拟合，architecture、data quality、optimizer、tokenizer 与 evaluation 改变后，系数和转折点都可能变化。
\end{warningbox}

\subsection{Outer loop 与 inner loop：参数学习和上下文适配}

训练阶段的 outer loop 用 SGD 更新参数：

\[
\theta_{k+1}=\theta_k-\eta\nabla_\theta \mathcal L(\theta_k;\mathcal B_k).
\]

推理时参数 $\theta$ 固定，prompt examples 改变每层 activations $H^{(\ell)}$，形成行为上的 inner-loop adaptation。称其为“在 activations 中学习”是有用描述；称其执行字面 gradient descent 则需要额外机制证据。

\lecturefigure{slide-56-in-context-learning.jpg}{In-context learning：示例数量增加时，无参数更新的任务表现改善。}{Stanford Online 官方视频 00:56:31--00:57:23。}

\begin{knowledgebox}{读图：区分三种变化}
Outer loop 改 weights，持续影响以后所有输入；in-context adaptation 改当前 forward 的 hidden states，只在 context 生命周期内生效；retrieval 或 tool use 还会改变外部输入。三者都能改善回答，却对应不同因果路径。把 prompt 效果都叫“fine-tuning”会掩盖系统行为。
\end{knowledgebox}

\subsection{Chain of thought：更多中间 token 也是更多计算步骤}

Chain-of-thought prompting 让模型在最终答案前生成中间推理文本。对自回归 Transformer，每个新 token 都触发一次新的 forward step，并把先前生成内容加入 context；因此中间 token 不只是解释，也提供额外 sequential compute 与可读 scratch space。但可读步骤仍可能包含错误，不能自动作为忠实内部因果链。

\lecturefigure{slide-57-chain-of-thought.jpg}{Chain-of-thought：用示例诱导中间推理步骤，提高部分任务表现。}{Stanford Online 官方视频 00:57:23--00:58:16。}

\begin{warningbox}{读图：正确答案与忠实解释是两项指标}
CoT 可能提高 arithmetic、symbolic 或 compositional tasks 的准确率，也可能产生事后合理化。评估时至少分开看 final answer、step validity、robustness to prompt perturbation 与 hidden reasoning fidelity。课堂用它支持“context 可配置计算”，并没有证明自然语言步骤等于模型内部真实算法。
\end{warningbox}

\subsection{三种性质必须同时成立}

Karpathy 最终把成功归纳为 expressive、optimizable、efficient。表达力让 attention/MLP 组合实现复杂 sequence function；可优化性来自 residual stream、normalization、较短梯度路径和成熟初始化；效率来自训练时在 token 维并行的大矩阵乘法。缺任一项，都可能得到理论强但训不动、能训但太慢，或高效却表达不足的模型。

\lecturefigure{slide-58-karpathy-tweets-transformer-history.jpg}{Karpathy 的高层总结：Transformer 同时优化表达力、可优化性与效率。}{Stanford Online 官方视频 00:58:16--00:58:45。}

\begin{knowledgebox}{术语消化：三项性质分别回答什么}
\begin{itemize}
  \item \term{expressive}：forward graph 能否表示复杂、数据依赖的计算；
  \item \term{optimizable}：gradient descent 能否稳定找到有用参数；
  \item \term{hardware-efficient}：主要算子能否转化为 GPU 擅长的大规模并行矩阵运算；
  \item \term{scalable}：前面三项在增大参数、数据与设备时是否仍保持可预测收益。
\end{itemize}
\end{knowledgebox}

\lecturefigure{slide-59-natural-language-model-history.jpg}{自然语言模型历史的重新解读：架构、规模与任务接口逐步汇合。}{Stanford Online 官方视频 00:58:45--00:59:17。}

这组高层推文是讲者面向大众的压缩叙事，不能替代论文细节。它的价值是把历史从“模型名字年表”改写为“计算图怎样变化”：固定窗口变成 recurrent state，单一 state 变成可寻址 memory，串行 recurrence 变成并行 attention，单任务 predictor 又在大规模 context 中表现出 runtime task adaptation。

\subsection{Shallow-wide 与 long-thin：硬件效率也是模型能力来源}

RNN 在时间上形成 long-thin graph，必须依次处理 token；Transformer 训练时形成 shallow-wide graph，同一层所有位置可并行，监督到输入的层级 hop 较少。Attention 的 $T^2$ 代价仍昂贵，但 dense matrix multiplication 极适合 GPU。更高设备利用率允许训练更大模型，而深度学习中规模本身会改变可达性能。

\lecturefigure{slide-60-general-purpose-computer.jpg}{GPT 作为 runtime-reconfigurable computer：prompt 是自然语言程序，completion 是执行轨迹。}{Stanford Online 官方视频 00:59:17--01:01:11。}

\begin{importantbox}{“通用计算机”是类比，不是无限能力证明}
固定 weights 的 GPT 可以因 prompt 而执行翻译、分类、抽取、格式转换或多步生成，类似同一硬件加载不同程序。但有限 context、有限 precision、训练分布、随机采样与错误累积都限制可靠性。类比强调 runtime reconfiguration，不保证任意程序可正确执行，也不保证输出满足形式规范。
\end{importantbox}

\teachervoice{Q\&A 中 Karpathy 进一步比较：RNN 即使原则上能实现任意程序，也可能因 long-thin serial graph 而难优化、难并行；Transformer 的 shallow-wide graph、residual pathways 与 LayerNorm 让监督更快流向输入。第三项“硬件效率”常被低估，却决定实际能训练多大。}

\subsection{本章小结}

Transformer 的成功来自 package-level synergy：attention/MLP 提供表达力，residual/normalization 提供可优化性，宽并行矩阵运算提供硬件效率，规模化 next-token training 又产生可由 prompt 配置的 in-context behavior。Scaling、few-shot 与 CoT 是行为证据；“inner optimizer”和“通用计算机”仍是需要边界的解释模型。

